\subsection{问题三分析}
系统长期动力学行为的识别比短期政策效应分析更依赖简洁的模型结构。若沿用前两问的高维划分，参数空间将急剧膨胀，反而难以归纳系统轨道类型。因此，我们在问题三中采用"基础资源供给-中层鱼类转移-顶级捕食者反馈"的最简三级食物链闭环，将环境承载力变化抽象为无量纲控制参数$K$，在同一参数轴上分析系统从规则有界振荡向复杂非线性波动的演化过程。

仅靠轨道形态不足以支撑系统行为分类，还需结合统一时间窗内的后期轨道统计量。通过综合分析波动幅度、峰谷差及峰间间隔的离散程度，才能为系统行为分类提供可重复验证的定量依据。

